\documentclass[runningheads]{llncs}
\usepackage{graphicx}
\usepackage{hyperref}
\usepackage{amsmath}
\usepackage{listings}
\usepackage{xcolor}

\title{Retrieval-Augmented Copilot for Quantum Simulation}
\titlerunning{RAG Copilot for Quantum Simulation}

\author{Your Name \and Team Member 2 \and Team Member 3}
\authorrunning{Your Name et al.}
\institute{Your University, Department\\
\email{youremail@uni.edu}}

\begin{document}
\maketitle

\begin{abstract}
We present a retrieval-augmented assistant specialized for quantum simulation that generates grounded explanations and runnable code in Qiskit, Cirq, and PennyLane. Our system combines dense retrieval over curated quantum corpora with a code-capable language model and a sandboxed execution loop. We evaluate faithfulness, code execution success, task accuracy (VQE/QAOA/tomography), and latency, demonstrating advantages over LLM-only baselines.
\keywords{Retrieval-Augmented Generation \and Quantum Computing \and Qiskit \and Cirq \and PennyLane}
\end{abstract}

\section{Introduction}
% Problem, motivation, contributions

\section{Related Work}
% RAG, vector search, code LLMs, quantum algorithms and SDKs

\section{Method}
% Ingestion, indexing, retrieval, prompting, code generation, execution, guardrails

\section{Experiments}
% Datasets, metrics, baselines, ablations

\section{Results}
% Tables/figures, analysis

\section{Discussion}
% Limitations, ethical/safety, future work

\section{Conclusion}
% Summary

\bibliographystyle{splncs04}
\bibliography{references}

\end{document}



